\subsection{F-01 --- \fillme{Short, specific finding title, e.g. "Certificate Pinning Disabled in Production Build"}}
\label{sec:f01}

\begin{center}
\begin{tabularx}{\linewidth}{@{}p{3.8cm} X p{3.8cm} X@{}}
\toprule
\rowcolor{TableHead}
\multicolumn{4}{l}{\textcolor{white}{\bfseries\sffamily Finding Metadata}} \\
\midrule
\textbf{ID}       & F-01 & \textbf{Severity}  & \critical \\
\rowcolor{TableAlt}
\textbf{Platform} & \bothplat & \textbf{Exploitability} & \fillme{Trivial / Moderate} \\
\textbf{Component} & \fillme{e.g. network layer / OkHttp client} & \textbf{Auth Required} & \fillme{None / Authenticated} \\
\bottomrule
\end{tabularx}
\end{center}

\subsubsection{Description}
\fillme{Describe the vulnerable component and how it was identified (static disassembly finding,
or dynamic Frida-hook observation).}

\subsubsection{Steps to Reproduce}
\begin{enumerate}
  \item \fillme{Step 1 -- e.g. install a Burp CA certificate and configure the device proxy.}
  \item \fillme{Step 2 -- e.g. launch the app and observe traffic is intercepted without a pinning error.}
\end{enumerate}

\subsubsection{Impact}
\begin{criticalbox}
\fillme{What an attacker in a privileged network position (public Wi-Fi, malicious VPN profile,
compromised router) can now do -- e.g. read/modify all API traffic including auth tokens.}
\end{criticalbox}

\subsubsection{Remediation}
\begin{itemize}
  \item \fillme{e.g. implement certificate/public-key pinning consistently across all build variants.}
  \item \fillme{e.g. fail closed (refuse the connection) rather than fail open when pinning validation errors.}
\end{itemize}
